\section{Introduction}
\label{sec:intro}

Cloud providers are moving latency-sensitive and regulated workloads onto confidential
virtual machines (CVMs): guest memory is transparently encrypted and integrity-protected by
the CPU, and the hypervisor -- along with any co-resident tenant -- is placed outside the
trust boundary. AMD SEV-SNP and Intel TDX both make this guarantee for entire VMs rather than
individual enclaves, and providers such as Nexa Cloud now offer CVM instance types as a
drop-in replacement for ordinary VMs, with no code changes required from the tenant.

Memory encryption closes the most direct attack: a curious or compromised hypervisor can no
longer read guest pages off the wire. It does nothing, however, for resources that must stay
\emph{shared} for the economics of multi-tenancy to work at all. The last-level cache (LLC) is
physically partitioned across sockets, not tenants, so a victim CVM and an attacker VM
scheduled on the same socket still contend for the same cache ways and observe each other's
occupancy through timing. Prime+Probe and Flush+Reload style attacks that were originally
demonstrated against ordinary co-resident VMs apply to this setting largely unmodified --
encryption protects the data the victim computes on, not the memory-access pattern it computes
with.

Cache partitioning defenses eliminate the channel by construction, but they do so by giving
up the very thing multi-tenancy is sold on: a fixed, static carve-up of the LLC caps the
number of tenants a socket can host regardless of their actual working sets, and providers
have been reluctant to ship it as a default. Detection is the alternative in production today,
but existing detectors were built for ordinary VMs and lean on signals a CVM host does not
have. Guest-side counters or page-fault instrumentation require modifying or observing the
tenant's kernel, which is precisely what SEV-SNP and TDX exist to prevent; the hypervisor
cannot install anything inside the CVM's trust boundary without invalidating its attestation.
A detector for this setting has to work from what the host can see \emph{outside} the guest --
and it has to be cheap enough to run continuously on every socket, not just when an incident
is already suspected.

This paper presents \sysname, a hypervisor-side detector that meets both constraints. \sysname{}
observes only per-core hardware performance counters -- LLC references, LLC misses, and
per-core eviction counts that the CPU already exposes to the host -- and never inspects the
victim CVM's memory, registers, or hypercalls. It turns counter deltas into a contention
feature between candidate attacker/victim core pairs, smooths the feature with a streaming
classifier, and on a sustained detection re-pins the suspected attacker's virtual core to an
isolated cache partition instead of migrating or killing the victim. The victim's schedule,
memory layout, and attestation are untouched throughout.

Our contributions are:
\begin{itemize}[leftmargin=1.4em]
  \item A guest-oblivious detection architecture for cross-VM cache contention attacks that
  requires no instrumentation, cooperation, or visibility inside a confidential VM
  (Section~\ref{sec:design}).
  \item A streaming classifier with an explicit recall/false-positive trade-off knob, evaluated
  against a reimplemented \baseline{} baseline and a naive fixed-threshold detector
  (Section~\ref{sec:evaluation}).
  \item An isolation response that contains a detected attacker without touching the victim's
  scheduling, memory, or attestation state (Section~\ref{sec:design}).
  \item An evaluation on three production-like host traces showing 94.6\% recall at 1.8\% FPR
  and 2.1\% victim throughput overhead, with a median 340\,ms time to isolate a sustained
  attacker (Section~\ref{sec:evaluation}).
\end{itemize}
